\section*{Appendix A --- Property definitions}
\label{app:properties}

\begin{itemize}
\item \textbf{Net charge @ pH 7.4}: Henderson--Hasselbalch over D/E/C/Y/H/K/R
  with standard pKa, plus termini, divided by length.
\item \textbf{pI}: bisection root of the charge curve.
\item \textbf{GRAVY}: Kyte--Doolittle mean.
\item \textbf{Surface apolar fraction}: exposure from the \emph{input}
  backbone (CA neighbours within 10\,\AA; the less-buried half of residues
  counts as exposed); fraction of \{G,P,A,V,I,L,M,F,W,Y\} at exposed sites.
  \emph{Validation}: against a Shrake--Rupley SASA computed on N/CA/C/O with a
  1.4\,\AA{} probe, the CA-burial exposure
  agrees on 65--76\% of residues (chance 50\%) and Spearman-correlates with
  real SASA at $\rho = -0.38$ to $-0.59$ on the
  six backbones tested. The proxy is therefore informative but coarse, and
  absolute surface numbers should not be
  compared with SASA-based hydrophobicity scores such as SAP.
\item \textbf{Aggregation 5-mer fraction}: fraction of 5-residue windows that
  are $\geq$80\% apolar.
\item \textbf{Recovery}: identity to the native sequence at aligned positions.
\item \textbf{Diversity}: mean pairwise Hamming distance over the designs of one
  (arm, backbone) cell (16 in the main run).
\end{itemize}
